\section{Experiments}

In this section, we describe the experimental setup, methodology, and results for evaluating the Trust-Enhanced Graph-Based Recommendation Model (GTERM). Our experiments are designed to rigorously test the reliability, transparency, and accuracy improvements of GTERM compared to baseline models.

\subsection{Experimental Setup}

\textbf{Datasets:} We conducted experiments on four diverse datasets: Yelp, Gowalla, ML-10M movie ratings, and Amazon book ratings. Each dataset provides a different perspective on user preferences and interaction dynamics, facilitating a robust evaluation of the proposed model.

\textbf{Parameter Specifications:} All models were trained using PyTorch Geometric with a uniform hyperparameter setting unless specified otherwise. The key parameters were as follows:
\begin{itemize}
    \item \textit{Learning Rate:} 0.01
    \item \textit{Batch Size:} 128
    \item \textit{GCN Layers:} 2
    \item \textit{GAT Heads:} 8
    \item \textit{Dropout Rate:} 0.5
\end{itemize}
Trust score balancing parameter \(\alpha\) was set initially to 0.5 and adjusted during hyperparameter tuning.

\textbf{Evaluation Metrics:} We evaluated the models using Recall@10, Precision@10, and NDCG@10. These metrics provide insights into both relevance and ranking quality of the recommendations.

\subsection{Methodology}

\textbf{Preprocessing and Graph Construction:} The preprocessing stage included cleaning and normalizing datasets using min-max scaling and Z-score standardization, ensuring cross-compatibility. Trust scores calculated from interaction metrics were incorporated as edge weights in the bipartite graph, effectively embedding trust at the core of the model's structure.

\textbf{Graph Neural Networks (GNN) Modeling:} We employed a hybrid architecture combining GCN and GAT layers to facilitate effective feature propagation and attention-based trust prioritization. Node features were initialized using historical interaction data and refined through GNN layers to capture intricate user-item dynamics.

\textbf{Composite Score Calculation and Recommendation Generation:} The final recommendation scores were derived by combining GNN-predicted interaction scores with calculated trust scores using a tunable \(\alpha\) parameter. This composite score balances interaction fidelity with trust insights, enhancing both accuracy and transparency.

\textbf{Implementation Details:} The implementation was structured in modular code fashion, ensuring each component could be debugged and optimized independently. We utilized PyTorch Geometric's DataLoader for batch processing, improving memory management and computational efficiency.

\subsection{Results and Analysis}

\textbf{Performance Comparison:} Tables \ref{tab:performance} and \ref{tab:ablation} summarize the recommendation accuracy of GTERM against benchmark models like Collaborative Filtering (CF) and standard GNN-based recommenders.



\begin{table}[h] 
\centering 
\label{tab:performance} 
\begin{tabular}{lccc} 
\hline Model & Recall@10 & Precision@10 & NDCG@10 \\ 
\hline 
CF & 0.045 & 0.023 & 0.031 \\ 
GCN & 0.057 & 0.028 & 0.036 \\ 
GTERM & \textbf{0.067} & \textbf{0.034} & \textbf{0.043} \\ 
\hline 
\end{tabular} 
\caption{Recommendation Performance of GTERM} 
\end{table}
\textbf{Statistical Significance:} The improvements in recall and precision metrics were statistically validated using paired t-tests, confirming that GTERM's enhancements over baseline models are significant (p < 0.05).

\textbf{Ablation Studies:} Table \ref{tab:ablation} illustrates the impact of each component's removal on the model's performance, underscoring the importance of trust integration and attention mechanisms.

\begin{table}[h]
\centering
\label{tab:ablation}
\begin{tabular}{lccc}
\hline
Configuration & Recall@10 & Precision@10 & NDCG@10 \\
\hline
GTERM without Trust & 0.061 & 0.030 & 0.039 \\
GTERM without GAT & 0.064 & 0.032 & 0.041 \\
Complete GTERM & \textbf{0.067} & \textbf{0.034} & \textbf{0.043} \\
\hline
\end{tabular}
\caption{Ablation Study Results}
\end{table}

\textbf{Error Analysis:} A detailed examination of erroneous recommendations revealed that trust scores significantly rectify misleading low-score interactions, reducing anomalies prevalent in models lacking trust incorporation.

\textbf{Conclusion and Future Work:} The GTERM model shows promising improvements in recommendation accuracy and transparency by integrating trust metrics. Future work will explore dynamic trust evaluation and enhance attention mechanisms to further boost performance.

